\subsection{其解为规则函数之原函数的微分方程}

回顾（II，第 54 页，定义 3）：定义在区间 \(\mathbf{u}\) 上的向量函数 \(\mathrm{I} \subset \mathbf{R}\) 称为规则函数，如果它在 I 的每个紧子集上都是阶梯函数的一致极限；等价条件是：在 I 的每个内点处函数 \(\mathbf{u}\) 都有右极限和左极限，并且当这两个端点属于 I 时，在 I 的左端点处有右极限、在 I 的右端点处有左极限（II，第 54 页，定理 3）。本章中我们限于讨论这样的微分方程 (1)：它的每个解都是 I 上某个规则函数的原函数。显然，若对于 I 到 H 中的每个连续映射 \(\mathbf{u}\)，函数 \(\mathbf{f}(t, \mathbf{u}(t))\) 都在 I 上规则，则上述条件便满足；下面的引理给出这方面的一个充分条件：

引理 1. 设 \(\mathbf{f}\) 是 \(\mathrm{I} \times \mathrm{H}\) 到 \(\mathrm{E}\) 中的映射，使得：对每个 \(\mathbf{f}_{\mathbf{x}}\)，用 \(\mathbf{x} \in \mathrm{H}\) 表示 \(t \mapsto \mathbf{f}(t, \mathbf{x})\) 到 \(\mathrm{I}\) 中的映射 \(\mathrm{E}\) 时，下列条件满足：\(1^{\circ}\) \(\mathbf{f}_{\mathbf{x}}\) 对每个 \(\mathrm{I}\) 在 \(\mathbf{x} \in \mathrm{H}\) 上规则；\(2^{\circ}\) 映射 \(\mathbf{x} \mapsto \mathbf{f}_{\mathbf{x}}\)（从 \(\mathrm{H}\) 到由 \(\mathcal{F}(\mathrm{I}, \mathrm{E})\) 到 \(\mathrm{I}\) 中的映射所成的集合 \(\mathrm{E}\) 中）在给 \(\mathcal{F}(\mathrm{I}, \mathrm{E})\) 赋予紧收敛拓扑（《一般拓扑学》，X，第 278 页）时是连续的。在这些条件下：

\(1^{\circ}\) 对 I 到 H 中的每个连续映射 \(\mathbf{u}\)，函数 \(t \mapsto \mathbf{f}(t, \mathbf{u}(t))\) 都在 I 上规则；更确切地说，这个函数在点 \(t_0 \in \mathbf{I}\) 处的右（相应地，左）极限等于函数 \(t \mapsto \mathbf{f}(t, \mathbf{u}(t_0))\) 在点 \(t_0\) 处的右（相应地，左）极限。

\(2^{\circ}\) 若 \((\mathbf{u}_n)\) 是 I 到 H 中的映射序列，在 I 的每个紧子集上都一致收敛于 I 到 H 中的连续函数 \(\mathbf{u}\)，则函数序列 \(t \mapsto \mathbf{f}(t, \mathbf{u}_n(t))\) 在 I 的每个紧子集上一致收敛于 \(\mathbf{f}(t, \mathbf{u}(t))\)。

\(1^{\circ}\) 设 \(\mathbf{c}\) 是 \(\mathbf{f}(t, \mathbf{u}(t_0))\) 在点 \(t_0\) 处的右极限；对每个 \(\varepsilon > 0\)，存在 \(t_0\) 在 I 中的紧邻域 V，使得对 \(\| \mathbf{f}(t, \mathbf{u}(t_0)) - \mathbf{c} \| \leqslant \varepsilon\) 与 \(t \in V\) 有 \(t > t_0\)。另一方面，存在 \(\delta > 0\)，使得关系式

\[
\mathbf {x} \in \mathrm{H}, \quad \| \mathbf {x} - \mathbf {u} (t _ {0}) \| \leqslant \delta
\]

蕴含对一切 \(\|\mathbf{f}(s,\mathbf{x})-\mathbf{f}(s,\mathbf{u}(t_{0}))\|\leqslant\varepsilon\) 有 \(s\in V\)；若 \(W\subset V\) 是 \(t_{0}\) 在 I 中的邻域，使得对每个 \(\|\mathbf{u}(t)-\mathbf{u}(t_{0})\|\leqslant\delta\) 有 \(t\in W\)，则对 \(\|\mathbf{f}(t,\mathbf{u}(t))-\mathbf{c}\|\leqslant2\varepsilon\) 与 \(t\in W\) 有 \(t>t_{0}\)，这就证明了 c 是 \(\mathbf{f}(t,\mathbf{u}(t))\) 在点 \(t_{0}\) 处的右极限。

\(2^{\circ}\) 设 \(\mathbf{K}\) 是 \(\mathrm{I}\) 的紧子集；\(\mathbf{u}\) 在 \(\mathrm{I}, \mathbf{u}(\mathbf{K})\) 中的前者上连续，故后者是 \(\mathrm{H}\) 的紧子集；对每个 \(\varepsilon > 0\) 与每个 \(\mathbf{x} \in \mathbf{u}(\mathbf{K})\)，存在数 \(\delta_{\mathbf{x}}\)，使得对每个 \(\mathbf{y} \in \mathrm{H}\)、\(\| \mathbf{y} - \mathbf{x} \| \leqslant \delta_{\mathbf{x}}\) 且每个 \(t \in \mathbf{K}\)，有 \(\| \mathbf{f}(t, \mathbf{y}) - \mathbf{f}(t, \mathbf{x}) \| \leqslant \varepsilon\)。存在有限多个点 \(\mathbf{x}_i \in \mathbf{u}(\mathbf{K})\)，使得以 \(\mathbf{x}_i\) 为中心、以 \(\frac{1}{2} \delta_{\mathbf{x}_i}\) 为半径的闭球构成 \(\mathbf{u}(\mathbf{K})\) 的一个覆盖；令 \(\delta = \operatorname{Min}\left(\delta_{\mathbf{x}_i}\right)\)。据假设，存在整数 \(n_0\)，使得当 \(n \geqslant n_0\) 时，对每个 \(\| \mathbf{u}_n(t) - \mathbf{u}(t) \| \leqslant \frac{1}{2} \delta\) 有 \(t \in \mathbf{K}\)。而现在，对每个 \(t \in \mathbf{K}\)，存在指标 \(i\) 使得

\[
\left\| \mathbf {u} (t) - \mathbf {x} _ {i} \right\| \leqslant \frac {1}{2} \delta_ {\mathbf {x} _ {i}};
\]

于是有 \(\|\mathbf{u}_{n}(t)-\mathbf{x}_{i}\|\leqslant\delta_{\mathbf{x}}\)，从而对每个 \(\|\mathbf{f}(t,\mathbf{u}_{n}(t))-\mathbf{f}(t,\mathbf{u}(t))\|\leqslant2\varepsilon\) 与每个 \(t\in K\) 有 \(n\geqslant n_{0}\)。

在本节余下部分，I 表示一个包含于 R 中且不退化为一点的区间，H 表示赋范空间 E 中所包含的一个非空开集，而 f 表示 \(I \times H\) 到 E 中、满足引理 1 诸假设的一个映射。

命题 1. 设 \(t_0\) 是 I 的一点，\(\mathbf{x}_0\) 是 H 的一点；连续函数 \(\mathbf{u}\) 为方程 (1) 在 I 上的解并取值 \(\mathbf{x}_0\) 于点 \(t_0\)，其充分必要条件是它满足关系式

\[
\mathbf {u} (t) = \mathbf {x} _ {0} + \int_ {t _ {0}} ^ {t} \mathbf {f} (s, \mathbf {u} (s)) d s\tag{6}
\]

对每个 \(t\in \mathrm{I}\)

的确，由引理 1，若 \(\mathbf{u}\) 是 (1) 在 I 上的解，则 \(\mathbf{f}(t, \mathbf{u}(t))\) 是规则的，故 (6) 的右端有定义并等于 \(\mathbf{u}(t)\)，对每个 \(t \in I\) 皆然。反之，若 \(\mathbf{u}\) 是满足 (6) 的连续函数，则由引理 1，\(\mathbf{f}(t, \mathbf{u}(t))\) 是规则的，从而 \(\mathbf{u}\) 除去 I 的一个可数子集中的点以外有等于 \(\mathbf{f}(t, \mathbf{u}(t))\) 的导数。

推论. 在 I 的每个不同于该区间左（相应地，右）端点的点处，(1) 在 I 上的每个解 u 都有左（相应地，右）导数，它等于 \(\mathbf{f}(t,\mathbf{u}(t))\) 在该点处的左（右）极限。

命题 2. 若 \(\mathbf{f}\) 是 \(\mathrm{I} \times \mathrm{H}\) 到 \(\mathrm{E}\) 中的连续映射，则 (1) 在 \(\mathrm{I}\) 上的每个解都是严格解。

的确，这样的解 \(\mathbf{u}\) 是连续函数 \(\mathbf{f}(t,\mathbf{u}(t))\) 的原函数（II，第 66 页，命题 3）。

此外我们指出，\(\mathbf{f}\) 上的连续函数 \(\mathrm{I} \times \mathrm{H}\) 满足引理 1 的条件（《一般拓扑学》，X，第 286 页，推论 3）。

往后我们任意选取 \(t_{0} \in I\) 与 \(x_{0} \in H\)，并考察是否存在 (1) 在 I 上（或在 \(t_{0}\) 在 I 中的某个邻域上）的解，它取值 \(x_{0}\) 于点 \(t_{0}\)（也就是 (6) 的解，二者是一回事）。

